\section{Strict-frame pooling: a direct quartic in 217 integers}
\label{sec:r21-integer}

The principal equation is the literal quartic
\[
 P_{217}(x_0,\ldots,x_{216})=\sum_{j=0}^{153}R_j(\boldsymbol x)^2=0,
 \qquad \boldsymbol x\in\mathbb Z^{217},
\]
where every quadratic residual and every collected coefficient is supplied
in the accompanying data. In ordinary ZFC, its integer nonsolvability is
equivalent to $\operatorname{Con}(\mathrm{ZFC})$, through the retained
numerical C222 and finite ZF/ZFC-refutation calibration. This is not a proof
of consistency, a new proof-assistant closure, or a global minimum claim.

\begin{center}
\begin{tabular}{lrrrr}
\toprule
Profile & Variables & Rows & Terms & Height\\
\midrule
Low height &217&154&25,949&4,576\\
Smaller support &217&154&14,615&16,866\\
Low-height companion &218&155&15,631&4,576\\
Sparse companion &218&155&8,817&16,866\\
\bottomrule
\end{tabular}
\end{center}
Every degree is exactly four and every residual coefficient has absolute
value at most 64. These four profiles are mutually nondominated in the
four displayed costs. The matched Strict Horizons parents have 235/172
and 236/173 variables/rows: each branch saves 18 of each. The distinct
incoming Pooled Powers 225-coordinate endpoint has 7,867 terms and
height 128,446, so it retains a support advantage; the new endpoints improve
its variable, residual and height costs, not every cost simultaneously.

\paragraph{What is new, and what is credited.}
The advance is a compatible composition and fully instantiated certificate,
not a priority claim for its component lemmas. The incoming Strict Horizons
source already has
\[
 t=16+h^2,\qquad F=t^t,\qquad R=F,\qquad q=4z+1,
\]
and strict local ranges using only the $q^t$ denominator.
Source 158, \emph{Pooled Powers}, supplies fixed-base pooling and
Pell-difference coordinates. Source 159, \emph{Quadratic Projection 234},
supplies common-radical projection and integer row reduction. Source 157,
\emph{Linear CRT}, already uses an odd-linear beta multiplier in a different
bootstrap. We prove compatibility with the unchanged strict frame and the
mod-four radix, including the preliminary negative-beta possibility.

\subsection{The unchanged frame supports both power pools}
Before interpreting powers, the retained signed frame, semantic anchor and
pooled orientation constraints give canonical signed Pell pairs and
$|A|>W^W$. A literal retained-row identity gives
\[
 W=2+\sum_{j=1}^{34}a_j^2.
\]
Every actual base, index and output of all seventeen retained requests is
one of these affine atoms. They include $t,F+t,x,y+1,F,B$, the nine CRT
representatives, $N$, the actual global derived label, $q,q+1,BH$, and the
fourteen pooled outputs. Hence every required square is $<W$, and $t\ge16$
gives $W\ge258$. In the dense 235 chart the identity is row 152 plus twice
row 4 plus row 169; in the sparse 236 chart it is row 153 plus row 4. None of
these rows is deleted. Actual $y$ is stored\_y$-1$; this offset is checked.

Here is the credited pooling argument, included to expose its hypotheses.
For a nonzero base $n$, let $X_i,Y_i$ be canonical signed Pell coordinates
at index $k_i$, put $Z_i=AY_i$, and assume $n^2,m_i^2,k_i^2<W$ for $r\le15$
requests. Write
\[
 T=2An-n^2-1,\quad D_i=m_i-X_i+Z_i,\quad
 H_n=\sum_iY_i(m_i-X_i),\quad S_n=\sum_i(D_i^2-Y_i^2).
\]
Two new coordinates suffice, with quadratic equations
\begin{equation}\label{eq:r21-power-pool}
 K_n-2H_n-(2A-n)L_n=0,\qquad S_n-nK_n+L_n=0.
\end{equation}
Indeed,
\[
 \sum_i(D_i-nY_i)^2=S_n-2nH_n-T\sum_iY_i^2,
\]
so $T$ divides the squared-error sum. Let
$b=\lfloor\sqrt{W-1}\rfloor\ge16$. Since $T\equiv-1\pmod n$, negative
Pell indices can be cleared by multiplication by $n^b$. Thus $T$ divides
\[
 E=\sum_i[m_i n^b-n^{b+k_i}]^2,
 \quad 0\le E\le4r b^{4b}<W^{2b+1}\le W^W<|T|.
\]
The last bound follows from
$|T|\ge2|A||n|-n^2-1\ge2|A|-W>W^W$.
Therefore $E=0$. At $|n|\ge2$, each $k_i\ge0$ and $m_i=n^{k_i}$;
a negative index would give the impossible integer identity
$m_i n^{-k_i}=1$. At $n=\pm1$ the signed unit output is exact without a
nonnegative-index claim. Conversely true powers give the divisibility
needed to set $L_n=(S_n-2nH_n)/T$ and then $K_n=2H_n+(2A-n)L_n$.

Apply this theorem to the eleven requests at $q+1$ and the three at $q$.
The initial requests $F=t^t,N=B^x,q=(BH)^x$ remain individual. The current
chart has just one old explicit modulus, $T_+$, because
\[
 T_q=T_+-2A+2q+1.
\]
Deleting fourteen quotients and that one modulus, and adding four pooled
coordinates, saves eleven variables and eleven residuals, not twelve.
All old power rows, retained $AY$ gates, norm rows and offsets are checked
literally before any deletion.

\subsection{The odd-linear beta bootstrap in the correct order}
Replace $v_\beta=(2w+1)^2$, $y=Fv_\beta$ by
\begin{equation}\label{eq:r21-linear-beta}
 y=F(2w+1).
\end{equation}
The monic multiplier contraction saves one variable and one residual.
Its multiplier may initially be negative, so the parent's positive-beta
argument cannot simply be quoted.

The retained power request first proves $F=t^t>1$. Consequently $y$ is a
nonzero multiple of $F$. Since $q=4z+1$, both $q\ne0$ and $|q+1|\ge2$
hold immediately. Interpret the $q+1$ pool first: it gives nonnegative
actual representatives, hence $H=2+\sum a_\nu\ge2$. The retained index
row $B\mu=C+F+1$ has positive right-hand side, so $B\ne0$ and $|BH|\ge2$.
The individual self-radix request gives $x\ge0$ and $q=(BH)^x$.

If $x=0$, the last endpoint
\[
 x=1+[1+(t+1)y]u_t
\]
would make $1+(t+1)y$ equal to $1$ or $-1$. This would give $y=0$ or
$(t+1)y=-2$, both impossible for a nonzero multiple of $F>1$ and $t\ge16$.
Thus $x\ge1$ and $|q|\ge2$. The $q$ pool now forces its requested index
$y+1\ge0$. Hence $y\ge-1$, and its nonzero-multiple property strengthens
this to $y\ge F>0$.

The original endpoint inequalities can now be used: from
$x=F+1+(y+1)u_0$ one gets $u_0\ge0$, then $x\ge F+1$, $u_t\ge1$, and
\[
 x\ge2+(t+1)y\ge2+(t+1)F.
\]
The inherited signed short-lift argument authentically recovers
$B=\binom{F+t}{t}>0$, and then $q=(BH)^x>0$. All strict ranges and local
rows follow as in Strict Horizons. The boundary $y=F$ is still excluded
\emph{after} local decoding: it would give $S_0=0$ and $S_t=1$, while
positive weighted updates propagate zero. Thus the source's label-product
and Horn-closure argument gives a finite C222 refutation.

\subsection{Completeness and unchanged arithmetic budgets}
The lcm-selector hypothesis remains valid:
$\Lambda_t\le4^t$ gives $F=t^t\ge16^t\ge\Lambda_t^2$, and the selector
requires a weak inequality. Its reverse direction covers every prime power
of $B$, not only row-isolating primes. The diagonal label-polynomial bound
is also unchanged. After equal nonunit factor counts $d\le t$ have been
recovered, each corresponding coefficient lies in $[1,F]$, so its
\emph{difference} has magnitude $<F$, even if an individual coefficient
is $F$. Least-degree reduction modulo $F$ recovers the whole polynomial.

For completeness choose an even $h$ in the cofinal family $t=16+h^2$,
strictly beyond the finite proof length and every local parameter and
local derived label. This is a witness choice, not a new restriction on
arbitrary integer solutions. Then $t$ is even and $v_2(F)\ge t>v_2(t!)$.
The parent constructs arbitrarily large valid $y=Fa^2$ with $a$ an odd
multiple of the odd part of $t!$. Equation~\eqref{eq:r21-linear-beta}
allows the same values, using $w=(a^2-1)/2$; an old square parameter maps
explicitly to $2w_{\rm old}(w_{\rm old}+1)$.
The odd beta moduli permit an arbitrarily large even $x$. The binomial is
odd, the prescribed representative parities make $H$ odd, and hence
$q=(BH)^x$ is an odd square, supplying $z=(q-1)/4$. The unchanged all-primes
CRT, short tails, norm and canonical Pell completion supply the other
witnesses. The converse of~\eqref{eq:r21-power-pool} supplies each new
$K_n,L_n$. This establishes the reverse implication.

\subsection{Exact maps, complete coefficients, and evidence}
Fourteen integral Pell-difference shears and a retained-norm row subtraction
reduce density. Six common-Hessian-radical directions then yield integral
quadratic zero-set bijections: along a direction $d$,
$R_i(u+sd)=R_i(u)+sL_i(d)$; a unit coordinate of $d$ and a unit row slope
solve one coordinate without increasing degree. The six directions,
restorations and inverse identities are frozen in certificates. Each final
row operation $R_i\leftarrow R_i+kR_j$ has the inverse with $-k$.

The semantic pooling/beta changes preserve \emph{zero existence}; they are
not advertised as bijections of every predecessor witness tuple. The later
coordinate and row transformations do have exact inverses on zero sets.
The independent checker imports no producer, checks the full frame,
pooling identities, all inverses, every collected quartic coefficient,
liveness, metrics and corruption rejection. Normal Python and Python
\texttt{-O} both pass. Reconstruction uses frozen certificates and a fresh
temporary directory, never an optimizer. The original arithmetic-parent
archive, proofs, notices and optional parent replay are retained.
Finite checks do not replace the unbounded mathematical proof or a
transitive proof-assistant build.
